\section{Introduction}
\label{sec:intro}

Practical video compression has long been dominated by spatial-temporal
hybrid codecs in the HEVC/VVC tradition, in which intra-coded keyframes
anchor a sequence of motion-compensated predictions and small per-block
residuals are entropy-coded. Learned video codecs of the
DCVC family~\cite{li2024neural,jia2025dcvcrt} have closed and in many cases
overtaken the rate--distortion (RD) frontier of these traditional schemes
while preserving the same Group-of-Pictures (GOP) intuition: a small,
strongly conditioned model produces a deterministic reconstruction, and
the transmitted bits describe how much the actual frames differ from that
reconstruction.

A complementary line of work has recently shown that pretrained
\emph{generative} video models can themselves serve as the decoder of an
ultra-low-bitrate codec. In this setting, the encoder transmits not a
residual against a deterministic prediction, but a compact specification of
which sample of the model's stochastic generation process to reproduce.
GVCC~\cite{zeng2026gvcc} instantiates this idea by converting a
rectified-flow video generator into a marginal-preserving stochastic process,
and replacing the per-step Gaussian noise of that process with codebook-
indexed innovations that the encoder selects to match the ground-truth video.
The result is a zero-shot codec---no training of the generator or any
auxiliary network---that on UVG achieves substantially lower LPIPS than
HEVC, VVC, and even strong learned codecs at sub-$0.01$~bpp.

The clear weakness of this design is decoder time. Each GOP requires
$T \approx 20$ forward passes through a multi-billion-parameter video
generator, plus a per-step codebook search; on the published
Wan~2.1-14B~\cite{wan2025} backbone, a single $33$-frame 1080p GOP takes
several hundred seconds to decode on a high-end GPU. This is one to two
orders of magnitude above the operating point of DCVC-RT and rules out the
codec in any latency-sensitive scenario.

In this paper we present \textbf{GVCCTurbo}, a turbo-charged variant of the
same generative-codec design. The central design constraint is
\emph{trajectory preservation}: the decoder should still traverse the same
fine-grained SDE/codebook correction process, because those per-step
innovations are the transmitted residual channel of the codec. We therefore do
not simply reduce the number of solver steps. Instead, we reduce the cost of
some steps while keeping the codebook trajectory reproducible. Our two
contributions are simple, trajectory-geometric, and \emph{require no
retraining}.

\paragraph{Predictive skip via clean-endpoint caching.}
The dominant per-step cost is the video generator's velocity evaluation.
We replace roughly half of these evaluations along the SDE trajectory with a
cached reconstruction. The naive way to cache---storing the previous step's
velocity and reusing it---silently drifts, because the velocity is a function
of both state and time. Instead, we cache the generator's
\emph{clean-latent estimate} (the trajectory endpoint) and reconstruct the
skipped step's velocity from the current state and current time. Under the
linear flow interpolant, the clean endpoint is invariant along the trajectory,
so this scheme is geometrically faithful within a local linearization. We show
empirically that the cache has an effective radius of one step: a fixed
alternating skip/run schedule, combined with a one-step warmup after each
resolution transition, is Pareto-optimal under this caching scheme; more
aggressive multi-step skipping is unstable.

\paragraph{Multi-resolution coding with a unified target.}
A natural way to cut decoder cost is to run the early high-noise iterations
of the generator at a lower spatial resolution and only switch to the
target resolution near the clean end---a video-coding analogue of progressive
spectral diffusion~\cite{xiao2026spd}. We find that the obvious instantiation
(encode at low-res VAE, switch to high-res VAE via a spectral lift) hits a
structural ceiling of roughly $1.4$~dB on UVG that persists across backbone
scale and stage count, and that does not close under a joint oracle. The
underlying cause is that the low-resolution VAE latent and the DCT-low
projection of the high-resolution VAE latent are \emph{different objects},
not aligned by the spectral lift used at transition. Unifying both stages
in the high-resolution latent domain---using the DCT-low projection of the
high-res latent as the low-stage codebook target---collapses the ceiling to
$0.2$~dB at the same bitrate, recovering nearly all of the gap with a one-
line change in the codec target.

\paragraph{Negative results as design constraints.}
Several plausible accelerations fail for principled reasons. Band-limiting
the codebook innovation breaks the full-spectrum noise distribution expected
by the SDE; caching velocity remembers the previous direction rather than the
current destination; multi-step cache reuse compounds its own extrapolation
error; and a naive low-resolution VAE target moves the codec into the wrong
latent domain. These failures shape the final design: GVCCTurbo reduces
generator cost only when the missing computation can be replaced without
changing the stochastic correction channel that carries the compressed bits.

\paragraph{GOP-level evaluation across conditioning modes.}
GVCC is evaluated in three GOP-conditioning modes: text-only (T2V), with
a single anchor frame (I2V) with autoregressive chaining, and with both
endpoint frames (FLF2V) with boundary sharing across consecutive GOPs.
We apply our acceleration mechanisms primarily to T2V and FLF2V; the
I2V mode is included by reusing the published GVCC-I2V configuration
without re-evaluation. Our main full-UVG measurement is at $720$p in
T2V mode: the full acceleration stack reduces per-GOP decoder
wall-clock to $19$~s with PSNR within $0.3$~dB of the single-resolution
baseline at the same bitrate. At $1080$p we report \emph{preliminary
half-UVG} results (first $9$ GOPs per sequence): in FLF2V mode the
turbo pipeline matches the bitrate of the published GVCC-FLF2V and
exceeds it by $+0.5$~dB PSNR and $-14\%$ LPIPS at approximately
$2\times$ the decode speed, and in T2V mode (Wan-T2V-1.3B) it reaches
$0.0030$~BPP at $108$~s per GOP. Full-UVG $1080$p measurements are the
next step.

The remainder of the paper is organized as follows. Section~\ref{sec:related}
positions our work against traditional, learned, and generative video
codecs. Section~\ref{sec:method} describes the pipeline and derives both
acceleration mechanisms, including the trajectory-geometric arguments and
falsified alternatives. Section~\ref{sec:experiments} reports the main RD
and decode-time comparison on UVG. Section~\ref{sec:ablations} isolates the
contribution of each component and validates the one-step-radius result for
predictive skipping.
